% Cover letter for ACM Computing Surveys (editor-facing; not part of the manuscript).
% Compile with:  pdflatex cover_letter.tex   ->  cover_letter.pdf
\documentclass[11pt]{article}
\usepackage[margin=0.8in]{geometry}
\usepackage{parskip}
\setlength{\parskip}{0.4em plus 1pt minus 1pt}
\linespread{0.96}
\pagestyle{empty}
\begin{document}

\noindent\textbf{Date:} July 2026\\
\textbf{To:} The Editor-in-Chief, ACM Computing Surveys\\
\textbf{Re:} Submission of \emph{``Weights or Skills? A Survey of Robot-Learning Techniques:
from Action-Predicting Weights to Robots that Write their Own Skills''}\\
\textbf{Article type:} Survey\\
\textbf{Corresponding author:} [Name, affiliation, email]

\smallskip
\noindent Dear Editor-in-Chief,

We submit \emph{``Weights or Skills?''} for consideration as a Survey in ACM Computing Surveys.
The manuscript organizes the fast-moving field of robot learning along a single original axis,
\textbf{what a system ships}: a generalist policy can ship frozen neural-network \textbf{weights}
that map observations directly to actions (the vision-language-action, or VLA, family), or it can
ship executable \textbf{skills} as code that the robot writes, debugs, and improves from its own
experience (the code-as-policy family). To our knowledge, no existing survey organizes the field
by this weights-versus-skills contrast.

This structure is needed because the field is fragmenting into two research cultures that rarely
cite each other, and recent surveys track only one pole at a time. Placing both poles on a common
axis makes them directly comparable on what they ship, where they are evaluated, and whether they
keep improving after deployment. On top of that axis the survey makes two contributions that, to
our knowledge, no prior survey provides: (i)~it arranges code-as-policy methods by \textbf{degree
of self-improvement}, giving operational definitions of the feedback, memory, and search
mechanisms with necessary-and-sufficient conditions for each rung, and it identifies the sparsely
populated cell in which execution feedback, persistent skill memory, and evolutionary search
combine into one open-ended loop, which only a few very recent systems occupy; and (ii)~it
connects the taxonomy to the emerging \textbf{robot skill economy}, in which skill marketplaces
distribute one-tap skills across robots, surfacing open problems in adaptation, cross-embodiment
portability, provenance, safety verification, and standardisation.

We have checked CSUR's recent volumes and the wider literature to position the work. The nearest
CSUR surveys address adjacent but distinct territory: robot navigation via foundation language
models (pipeline-stage scope, navigation only), embodied intelligence as a synergy of morphology,
action, perception, and learning (morphology-centric), world models, and an earlier comprehensive
survey of hierarchical reinforcement learning that predates the LLM code-as-policy shift. Beyond
CSUR, the strongest incumbents each cover a single pole or a single sub-area: the
vision-language-action survey treats only the weights pole; the recent self-evolving-agents
surveys treat the self-improvement pole but are embodiment-agnostic rather than robot-specific; and
surveys of LLM-enhanced reinforcement learning, cross-domain policy transfer, and generative
robotic manipulation each own one of our supporting sub-sections. We cite these as the authorities
for their sub-areas and frame our value as the cross-cutting synthesis plus the two novel
contributions above, not as fresh coverage of any one topic. We found no survey, in CSUR or
elsewhere, that organizes robot learning by weights versus skills, arranges code-as-policy by
degree of self-improvement, or treats the robot skill economy.

The treatment is comprehensive. It draws on 311 references and 302 systems: 77 placed in the
taxonomy and compared across six per-family tables, plus 225 further works catalogued in a
landscape appendix, spanning January 2016 to July 2026 and assembled by a documented two-pass,
PRISMA-style methodology with explicit inclusion and core-versus-appendix placement criteria. The
package includes an original taxonomy figure in the introduction, a comparison table that
organizes the surveyed literature, a family-by-six-axes spec matrix, and an open-problems agenda
whose future directions are anchored on measurable quantities and named testbeds.

The manuscript is 40 pages in the \texttt{acmsmall} format and is prepared double-anonymous for
review. The work is original, has not been published previously, and is not under consideration
elsewhere. The authors declare no conflicts of interest [confirm before sending]. We suggest the
following reviewers, chosen for breadth across both poles and free of collaboration conflicts:
[list 4+ names, affiliations, emails]. Thank you for considering the manuscript.

\enlargethispage{2\baselineskip}
\smallskip
\noindent Sincerely,\\[2pt]
[Authors, on behalf of all co-authors]

\end{document}
